\subsection{The actual Case II residual links}
\label{subsec:sr-four-case-II-residual-links}

This subsection supplies the residual-path producer in Case II.
It starts with a residual point in the good stratum.
Producing that point from an arbitrary original
\(\mathrm{Star}\)-representation is the separate lattice step.

\begin{srlemma}[A strict class with a prescribed ramification line]
\label{lem:sr-four-ramified-line-class}
Let \(F\) be a number field, \(\ell\ne2\), and \(M\) a finite
\(\mathbb F_\ell[G_F]\)-module. Fix \(m_0\ne0\) in \(M\)
and finitely many local groups at which the class must vanish
in \(H^1\). After allowing ramification at finitely many fresh
rational primes splitting completely in any prescribed finite
Galois extension containing \(F\), the field of \(M\), and
\(\mu_\ell\), there is a continuous cocycle \(c\) such that:
it is locally a coboundary at the fixed groups;
at every fresh inertia group its values lie in
\(\mathbb F_\ell m_0\); and at one such inertia element \(t\),
\[
                    \rho(t)=1,\qquad c(t)\ne0.             \tag{R4L1}
\]
It is unramified outside the original and fresh sets.
\end{srlemma}
\begin{proof}
Choose an original finite set \(S_0\) containing the required
local groups and all ramification of \(M\).
At \(S_0\) impose the zero local condition.
At a fresh prime with \(\rho|_{G_v}=1\) and
\(\mathrm Nv\equiv1\pmod\ell\), impose the condition
that the inertia value belongs to \(\mathbb F_\ell m_0\).
Local tame cohomology identifies this subspace with the
direct sum of \(M\), for its unramified value, and the
one-dimensional prescribed inertia line. Its dimension is
\(\dim_{\mathbb F_\ell}M+1\).
Its contribution
\(\dim L_v-\dim H^0(F_v,M)\) to the
Euler--Poitou--Tate formula is therefore \(1\).

For each fresh rational prime choose one of its primes in \(F\)
for this condition; at the other primes above it impose the
unramified condition, whose contribution is zero.
For \(h\) fresh rational primes the exact Selmer formula gives
\[
 \dim H^1_{\mathcal L_h}(F,M)
       -\dim H^1_{\mathcal L_h^\perp}(F,M^\vee(1))
                         =h+C,                            \tag{R4L2}
\]
where \(C\) depends only on the fixed original local data.
The space of classes unramified at all the fresh primes is
the fixed finite-dimensional Selmer space for \(S_0\).
Choose \(h\) larger than its dimension minus \(C\).
Since the dual dimension is nonnegative, (R4L2) gives a
class ramified at one fresh prime.
All cocycle representatives have the same inertia values
there, because the local action is trivial and local
coboundaries have zero inertia values.
This proves (R4L1).
The local conditions at \(S_0\) give the stated coboundaries.
Chebotarev supplies as many fresh rational primes with all
the required splitting conditions as needed.
This is an existence argument; no numerical bound involving
these auxiliary degrees will enter the final estimate.
\end{proof}

\begin{srlemma}[Simultaneous actual cup primitives]
\label{lem:sr-four-simultaneous-cup}
Let \(\alpha,\beta:G_F\to k^\times\) be the finite residual
root characters, with \(k\) finite of characteristic \(\ell\).
Given finitely many normalized \(\beta\)-cocycles \(y_j\),
there are, after the fresh ramification just described, an
\(\alpha\)-cocycle \(x\), its local splitting parameters \(m_v\),
and functions \(z_j\) with local parameters \(o_{jv}\), such that
\[
 \begin{split}
 x(gh)&=x(g)+\alpha(g)x(h),\\
 z_j(gh)&=z_j(g)+\alpha(g)\beta(g)z_j(h)
                         +y_j(g)\alpha(g)x(h),\\
 x|_{G_v}&=m_v-\alpha m_v,\\
 z_j|_{G_v}&=o_{jv}-\alpha\beta o_{jv}
                              -y_j\alpha m_v.             \tag{R4L3}
 \end{split}
\]
Moreover \(x\) is nonzero on \(\ker\alpha\).
If \(\alpha(g_0)\ne1\), \(\alpha(g_0)\beta(g_0)\ne1\),
and \(y_j(g_0)=0\), one may require
\(x(g_0)=z_j(g_0)=0\).
\end{srlemma}
\begin{proof}
Use the finite module
\[
 M=k\oplus k^J,\qquad
 \rho(g)(a,(b_j))=
 \bigl(\alpha(g)a,\
       (\alpha(g)\beta(g)b_j+\alpha(g)y_j(g)a)_j\bigr).
\]
The cocycle equation for \(y_j\) proves that \(\rho\) is a
representation.
Apply Lemma~\ref{lem:sr-four-ramified-line-class} with
\(m_0=(1,0,\ldots,0)\).
The first and remaining coordinates of its cocycle are
\(x,z_j\). The cocycle and local-coboundary equations give
exactly (R4L3), with the negatives of the corresponding
coordinates of the local correcting vectors.
At the nonzero fresh inertia element, all \(z_j\) are zero
and \(x\ne0\), while \(\alpha=1\).

For normalization put
\[
 b=\frac{x(g_0)}{1-\alpha(g_0)},\qquad
 a_j=\frac{z_j(g_0)}{1-\alpha(g_0)\beta(g_0)}.
\]
Replace
\[
 \begin{split}
 x(g)&\longmapsto x(g)-(1-\alpha(g))b,\\
 z_j(g)&\longmapsto z_j(g)
        -(1-\alpha(g)\beta(g))a_j+y_j(g)\alpha(g)b,\\
 m_v&\longmapsto m_v-b,\qquad
 o_{jv}\longmapsto o_{jv}-a_j .
 \end{split}
\]
Substitution proves every equation in (R4L3).
The new values at \(g_0\) are zero.
The value of \(x\) on \(\ker\alpha\), including its nonzero
fresh inertia value, is unchanged.
\end{proof}

We apply this with
\(\alpha=\bar\chi_1/\bar\chi_0\) and
\(\beta=\bar\chi_2/\bar\chi_0\).
Here \(\beta\) is the unramified quadratic root character of
Case II and is \(1\) on every ordinary local group.
The distinguished regular \(g_0\) makes
\(\alpha(g_0),\beta(g_0),\alpha(g_0)\beta(g_0)\ne1\).

\paragraph{The matrices and flags of the intermediate point.}
Set
\[
 U(x,y,t)=
 \begin{pmatrix}
 1&0&0&0\\x&1&0&0\\y&0&1&0\\t&y&-x&1
 \end{pmatrix}.
\]
Direct multiplication gives
\[
 U(x,y,t)U(x',y',t')
       =U(x+x',y+y',t+t'+yx'-xy'),\qquad
                         U(x,y,t)^\#=U(-x,-y,-t).         \tag{R4L4}
\]
In particular these are lower unipotent symplectic matrices.

Let \(p\) be the given good residual point, supported in
\[
 N_P=\{\beta_{\rm root},-\alpha_{\rm root},
                       -\theta_{\rm root},-\psi_{\rm root}\},
\]
with its lower flags supported in
\(Q=\{-\alpha_{\rm root},-\psi_{\rm root},-\theta_{\rm root}\}\).
Its \(20\) entry is \(y(g)\bar\chi_0(g)\), where
\[
 y(gh)=y(g)+\beta(g)y(h),\quad y(g_0)=0,\quad
 y|_{G_v}=n_v-\beta n_v,\quad n_v=p.n_v{}_{20}.             \tag{R4L5}
\]
These identities follow from the \(20\) entry of multiplication
and of the ordinary flag equation; no other entry contributes
in the support \(N_P\).

Use Lemma~\ref{lem:sr-four-simultaneous-cup} for this \(y\),
and write \(D(g)=\operatorname{diag}(\bar\chi_i(g))\).
Then
\[
 q(g)=U(x(g),y(g),2z(g)-x(g)y(g))D(g),\qquad
 q.n_v=U(m_v,n_v,2o_v-m_vn_v)                              \tag{R4L6}
\]
is the actual intermediate residual point.
Indeed the fixed pairing gives
\(D U(x,y,t)D^{-1}=U(\alpha x,\beta y,\alpha\beta t)\).
The second equation in (R4L3) makes the central coordinate
of (R4L6) satisfy exactly (R4L4).
The local equations make
\(q(g)=q.n_vD(g)(q.n_v)^{-1}\) on each ordinary group.
Thus all local flags and their ordered diagonal characters
are correct.
The normalization at \(g_0\), multiplier, and killed-inertia
conditions hold by construction.
Its support is \(Q\).
It has the same \(20\) entry and same local \(20\) flag
coordinate as \(p\), and a nonzero \(10\) entry.
These are the exact \(\mathrm{SamePs}\) and
\(\mathrm{Ent}(-\alpha_{\rm root})\) assertions.

\paragraph{Kill the introduced ramification before taking depth.}
The fresh primes may be chosen to split completely also in the
original residual point's field.
At a fresh inertia group, (R4L1) makes \(y=z=0\);
there (R4L6) has the form
\(U(x,0,0)\), of exponent \(\ell\), including when \(\ell=3\).
A fresh prime \(w\) satisfies \(w\equiv1\pmod\ell\).
The unique degree-\(\ell\) subfield of
\(\mathbb Q(\zeta_w)\) is real, unramified away from \(w\),
and totally tamely ramified at \(w\).
Its inertia subgroup is the kernel of the order-\(\ell\)
tame quotient. Consequently restriction to its compositum
with the current real field kills this introduced inertia.
For finitely many fresh primes take the compositum of these
real cyclic fields. It is abelian of \(\ell\)-power degree,
unramified at \(\ell\).
The original point was already unramified in the required
sense at these primes and remains so.

The nonzero \(x\)-entry survives every subsequent
\(\ell\)-power base extension.
In fact a cocycle for a nontrivial prime-to-\(\ell\)
character cannot vanish on a normal \(\ell\)-power-index
subgroup: that character is still nontrivial on the subgroup;
the cocycle equation and its conjugates would then force
the cocycle to vanish on the whole group.
The same argument applies to a nonzero normalized Levi
cocycle in (R4L5).
The distinguished element may be replaced by an
\(\ell\)-power of itself; its residual eigenvalues remain
pairwise distinct.

Restriction of flags is actual restriction, followed by
lower/upper Gaussian factorization.
For support \(N_P\), the first three principal minors of
\(r(g)n_v\) are products of the nonzero prescribed diagonal
characters. Thus this factorization exists.
Its new \(20\) flag coordinate is determined by
\[
       n'_{20}\bar\chi_0(g)
             =r(g)_{20}+\bar\chi_2(g)n_{20}.                \tag{R4L7}
\]
In particular two points with the same \(20\) data retain
that equality after restriction.
For lower support \(N_M\), the same argument uses the
\(10\) coordinate.
This is the required flag transport in the two
\(\mathrm{Same}\) assertions.

We now include these finitely many real cyclic fields in the
fixed abelian base used for the Ferrero--Washington bound,
and only afterwards choose the cyclotomic \(\ell\)-power
depth. Thus the rank-four bounds hold with one finite \(C\)
and \(d\ge30(C+1)^2\).
No quantitative assertion uniform in this auxiliary \(C\)
is needed.
The normalized induced seed exists at that final depth by
Proposition~\ref{prop:sr-induced-case-two-seed}.

\paragraph{The three polynomial paths.}
All statements below are over one finite coefficient field.
Let \(s\) be that seed's residual point.
It has support
\(V=\{\beta_{\rm root},-\alpha_{\rm root},
-\theta_{\rm root}\}\), lower flags in \(Q\), and
\(s.n_v{}_{20}\ne0\) at every ordinary place.

First join \(p\) to \(q\) by the affine path in the
parabolic with Levi root \(-\psi_{\rm root}\), keeping
their shared \(20\) data fixed.
This is a polynomial residual path.
One explicit check orders the two blocks as \((0,2)\)
and \((1,3)\). Its matrices have the form
\(\left(\begin{smallmatrix}A&0\\C&B\end{smallmatrix}\right)\),
where the two Levi blocks are fixed and \(C\) varies affinely.
Products of two differences in this off-diagonal block are
zero. Thus the representation and symplectic equations
are linear in the varying block.
The same block calculation for the lower flags proves the
ordinary equations along the whole path.

Next let \(q_0(g)=U(x(g),0,0)D(g)\).
Choose its lower flags
\[
        q_0.n_v=
          U(m_v,s.n_v{}_{20},-m_v\,s.n_v{}_{20}).           \tag{R4L8}
\]
Since \(\beta|_{G_v}=1\), (R4L4) shows that these diagonalize
\(q_0|_{G_v}\).
Thus \(q_0\) has the same \(10\) data as \(q\), and the same
\(20\) data as \(s\).
Join \(q\) to \(q_0\) by the corresponding affine
parabolic path with Levi root \(-\alpha_{\rm root}\).
Its radical consists of the other three negative roots
and is abelian; the identical two-block argument proves
all its equations.
Finally join \(q_0\) to \(s\) by the first kind of path,
keeping \(20\) data fixed.

The endpoints of these paths are good points.
The initial point \(p\) is good by assumption.
If its \(20\) entry is nonzero, \(q\) has both that entry
and the newly nonzero \(10\) entry.
If it is zero, the partially split good alternative has
nonzero \(20\) flag coordinates, retained by (R4L7);
then \(q\) lies in \(V\) and has its nonzero \(10\) entry.
The point \(q_0\) has nonzero \(10\) entry and every
\(20\) flag coordinate in (R4L8) is nonzero.
The seed is good by its explicitly proved residual shape:
either its \(10\) entry is nonzero, or its
\(\beta_{\rm root}\) and \(-\theta_{\rm root}\) entries
are nonzero and the required local root character is
nontrivial. That local nontriviality is the tower's
\(\mathrm{Two.bigL}\) condition.

The line property of the actual common space in
Section~\ref{subsec:sr-four-common-space-construction}
therefore realizes
\[
                         p\longleftrightarrow q
                          \longleftrightarrow q_0
                          \longleftrightarrow s            \tag{R4L9}
\]
by its actual polynomial residual maps.
Each consecutive pair has a common irreducible line
specialization beneath it.
The component propagation of
Proposition~\ref{prop:sr-four-core}, the actual modular ideal
of Section~\ref{subsec:sr-case-II-actual-pm}, and seed local
zero consequently propagate support along this finite chain.

All choices preceding (R4L9) are finite:
the locally constant residual coordinates lie in a finite
field; the auxiliary prime set and cyclic composita are
finite; and the final depth and seed have finite coefficient
fields. Their degree and the Ferrero--Washington constant
are used only to produce support. The final fixed-field
spectrum count is taken after descent and has no dependence
on any of them.

This reconstructs the Case II path producer
\srcid{Chain2.deep\_tower\_residue\_chain\_exists}.
The finite-module argument expands
\srcid{RamifiedLineCharacters.CupRaw}, \srcid{CupNorm},
and \srcid{GysinPsOf}; (R4L6) is the matrix construction in
\srcid{CompactIntegralFields.GysinAl}.
The restriction and flag transport are
\srcid{DiesEx}, \srcid{RestrT}, and \srcid{FlagFacts}.
The good residual lattice preceding this chain is constructed in
Subsection~\ref{subsec:sr-four-good-residual-lattice}.
The final block-extension and arithmetic endpoint assertions remain
separate from these residual paths.
